What happens after the last theorem of a book? This final tome of the series is a workshop in mathematical research,
and every tool in it is used on real mathematics, proved in full. It reads research papers the way mathematicians do:
the seven-shuffles theorem of Bayer and Diaconis, Selberg's central limit theorem for the zeta function, Wood's random
groups, and Harper's discovery that random multiplicative functions cancel better than chance allows. It discovers
formulas with a computer and lattice reduction, predicts the solutions of sums of three cubes with a heuristic, checks
proofs with the Curry--Howard correspondence behind proof assistants, and proves Jordan's theorem for polygons as a model
of careful writing. It ends with seven open problems, each with a partial result proved, and five projects that end
with a question nobody has answered. Each chapter opens with the story behind its ideas, ends with exercises and
complete solutions, and comes with a Python notebook.
